\documentclass{article}
\usepackage[T1]{fontenc}
\usepackage{lmodern}
\usepackage{graphicx}
\usepackage{amsmath,amssymb}
\usepackage[a4paper,margin=2cm]{geometry}
\usepackage[absolute,overlay]{textpos}
\setlength{\TPHorizModule}{1mm}
\setlength{\TPVertModule}{1mm}
\usepackage[indonesian]{babel}
\usepackage{hyperref}
\setlength{\emergencystretch}{2em}
% unit: Halaman 13

\begin{document}

% === Halaman 13 (llm) ===

\begin{itemize}
    \item Secara umum, jika pergeseran huruf sejauh $k$, maka:
    
    \begin{align*}
    \text{Enkripsi: } c &= E(p) = (p + k) \pmod{26} \\
    \text{Dekripsi: } p &= D(c) = (c - k) \pmod{26} \\
    k &= \text{kunci rahasia}
    \end{align*}
    
    \item Untuk alfabet berupa 256 karakter ASCII, maka:
    
    \begin{align*}
    \text{Enkripsi: } c &= E(p) = (p + k) \pmod{256} \\
    \text{Dekripsi: } p &= D(c_i) = (c - k) \pmod{256} \\
    k &= \text{kunci rahasia}
    \end{align*}
\end{itemize}

\end{document}
